\chapter{Les réactions allergiques}

\noindent\textit{Le système immunitaire, étudié au chapitre 7, nous protège
normalement des microbes. Mais il peut parfois réagir de façon excessive face à des
substances pourtant inoffensives : c'est l'allergie, un dérèglement de la réponse
immunitaire.}
\vspace{8pt}

\connecteBanner

\section{Qu'est-ce qu'une allergie ?}

\begin{definition}[Allergène]
Un \textbf{allergène} est une substance normalement \textbf{inoffensive} (pollen,
poussière, certains aliments, venin d'insecte, certains médicaments\dots) qui
déclenche, chez certaines personnes, une réaction immunitaire excessive.
\end{definition}

\begin{definition}[Allergie]
Une \textbf{allergie} est une réaction \textbf{exagérée} et \textbf{inappropriée} du
système immunitaire face à un allergène, alors que celui-ci ne représente aucun
danger réel pour l'organisme.
\end{definition}

\begin{exemple}
Le pollen de certaines plantes est totalement inoffensif pour la plupart des
personnes, mais déclenche des éternuements et un écoulement nasal chez les personnes
allergiques (rhume des foins).
\end{exemple}

\section{Le mécanisme de la réaction allergique}

\begin{propriete}[Deux étapes]
\begin{itemize}
\item \textbf{Premier contact — sensibilisation} : lors du premier contact avec
l'allergène, l'organisme le mémorise, un peu comme pour un microbe (voir chapitre 7),
sans provoquer de symptôme visible.
\item \textbf{Contact suivant — réaction allergique} : lors d'un contact ultérieur
avec le même allergène, certaines cellules libèrent en grande quantité une substance,
l'\textbf{histamine}, responsable des symptômes de l'allergie.
\end{itemize}
\end{propriete}

\begin{illus}
\begin{tikzpicture}[scale=0.85]
\node[font=\bfseries\small,align=center] at (0,1.3) {1\textsuperscript{er} contact};
\node[font=\scriptsize,align=center] at (0,0.55) {sensibilisation\\(sans symptôme)};
\draw[onbo,fill=onbo,fill opacity=0.12,line width=1.2pt] (0,-0.6) circle (0.7);
\node[font=\scriptsize] at (0,-0.6) {allergène};
\draw[->,onbmuted,line width=1.1pt] (1.1,-0.3) -- (2.9,-0.3);
\node[font=\bfseries\small,align=center] at (4.5,1.3) {2\textsuperscript{e} contact};
\node[font=\scriptsize,align=center] at (4.5,0.55) {libération\\d'histamine};
\draw[onbink,fill=onbink,fill opacity=0.15,line width=1.2pt] (4.5,-0.6) circle (0.85);
\node[font=\scriptsize,align=center] at (4.5,-0.6) {symptômes};
\end{tikzpicture}
\fig{Figure — Sensibilisation au premier contact, puis réaction allergique
(histamine) lors d'un contact ultérieur.}
\end{illus}

\begin{propriete}[Symptômes courants]
Éternuements, écoulement nasal, démangeaisons, éruption cutanée (urticaire),
gonflement, difficultés respiratoires. Dans les cas les plus graves, une allergie peut
provoquer un \textbf{choc anaphylactique}, une réaction générale brutale mettant la
vie en danger.
\end{propriete}

\begin{remarque}
Piège fréquent : l'allergène lui-même n'est pas dangereux ; ce sont les symptômes
provoqués par la \textbf{réponse du système immunitaire} qui peuvent l'être. C'est la
réaction de l'organisme, et non la substance, qui pose problème.
\end{remarque}

\section{Prévenir et traiter les allergies}

\begin{propriete}[Moyens de prise en charge]
\begin{itemize}
\item \textbf{éviction de l'allergène} : éviter tout contact avec la substance
identifiée comme responsable ;
\item les \textbf{antihistaminiques}, médicaments qui atténuent les effets de
l'histamine et soulagent les symptômes légers à modérés ;
\item en cas de choc anaphylactique, une injection d'\textbf{adrénaline} en urgence,
qui contre rapidement les effets dangereux de la réaction allergique.
\end{itemize}
\end{propriete}

\begin{exemple}
Une personne allergique aux arachides doit vérifier la composition des aliments
qu'elle consomme pour éviter tout contact avec cet allergène, seul moyen fiable de
prévenir une réaction allergique grave.
\end{exemple}

% ═══════════════════════════════════════════════════════════════
\section{L'essentiel à retenir}

\begin{synthese}
\textbf{Allergène.} Substance normalement inoffensive qui déclenche une réaction
immunitaire excessive chez certaines personnes.\\[4pt]
\textbf{Allergie.} Réaction exagérée et inappropriée du système immunitaire.\\[4pt]
\textbf{Mécanisme.} Sensibilisation lors du premier contact, puis libération
d'histamine lors des contacts suivants.\\[4pt]
\textbf{Choc anaphylactique.} Forme la plus grave d'allergie, mettant la vie en
danger.\\[4pt]
\textbf{Prise en charge.} Éviction de l'allergène, antihistaminiques, adrénaline en
urgence.\\[4pt]
\textbf{Piège fréquent.} L'allergène n'est pas dangereux en soi ; c'est la réaction du
système immunitaire qui produit les symptômes.
\end{synthese}

% ═══════════════════════════════════════════════════════════════
\section{Méthodes \& savoir-faire}

\methode{Reconnaître une réaction allergique}
Vérifier si les symptômes apparaissent après un contact avec une substance
habituellement inoffensive, et non lors d'un premier contact.
\begin{exemple}
Des éternuements systématiques au printemps, en présence de pollen, évoquent une
allergie plutôt qu'une infection.
\end{exemple}

% ═══════════════════════════════════════════════════════════════
\section{Exercices d'application}
\begin{multicols}{2}

\rubrique{Comprendre l'allergie}
\exo{}\diff{1} Qu'est-ce qu'un allergène ?

\exo{}\diff{2} Pourquoi dit-on que l'allergie est une réaction « inappropriée » ?

\rubrique{Mécanisme}
\exo{}\diff{2} Que se passe-t-il lors du premier contact avec un allergène ?

\exo{}\diff{2} Quelle substance est libérée lors d'une réaction allergique ?

\rubrique{Symptômes et prise en charge}
\exo{}\diff{1} Citer deux symptômes courants d'une allergie.

\exo{}\diff{2} Qu'est-ce qu'un choc anaphylactique ?

\exo{}\diff{1} Citer un moyen de prévenir une réaction allergique.

% ═══════════════════════════════════════════════════════════════
\end{multicols}

\section{Exercices d'entraînement}
\begin{multicols}{2}

\rubrique{Comprendre l'allergie}
\exo{}\diff{3} Expliquer pourquoi une personne peut être en contact avec un allergène
plusieurs fois avant de développer des symptômes.

\rubrique{Mécanisme}
\exo{}\diff{3} Comparer, en une phrase chacune, la réaction du système immunitaire
face à un microbe et face à un allergène.

\rubrique{Symptômes et prise en charge}
\exo{}\diff{2} Pourquoi les antihistaminiques soulagent-ils les symptômes de
l'allergie ?

\exo{}\diff{3} Pourquoi l'adrénaline est-elle administrée en urgence lors d'un choc
anaphylactique ?

\exo{}\diff{3} Une personne allergique aux piqûres de guêpe se fait piquer pour la
première fois sans réaction grave, puis fait un choc anaphylactique lors d'une
deuxième piqûre. Expliquer cette différence.

% ═══════════════════════════════════════════════════════════════
\end{multicols}

\section{Sujets type BEPC}

\sujetbac[Comprendre les allergies]
\begin{enumerate}[label=\arabic*)]
\item Définis les mots « allergène » et « allergie ».
\item Décris les deux étapes du mécanisme d'une réaction allergique.
\item Cite trois symptômes possibles d'une allergie.
\item Cite deux moyens de prévenir ou de traiter une réaction allergique.
\end{enumerate}

% ═══════════════════════════════════════════════════════════════
\section{Corrigés}
\begin{multicols}{2}

\corrige{de l'exercice 1}
Une substance normalement inoffensive qui déclenche une réaction immunitaire
excessive chez certaines personnes.

\corrige{de l'exercice 2}
Car elle se déclenche face à une substance qui ne représente en réalité aucun danger
pour l'organisme.

\corrige{de l'exercice 3}
L'organisme mémorise l'allergène (sensibilisation), sans provoquer de symptôme
visible.

\corrige{de l'exercice 4}
L'histamine.

\corrige{de l'exercice 5}
Par exemple : éternuements et démangeaisons.

\corrige{de l'exercice 6}
Une réaction allergique générale et brutale, qui met la vie en danger.

\corrige{de l'exercice 7}
Par exemple : éviter tout contact avec l'allergène identifié.

\corrige{de l'exercice 8}
Parce que le premier contact (ou les premiers contacts) ne provoque qu'une
sensibilisation silencieuse, sans symptôme ; les symptômes n'apparaissent qu'aux
contacts suivants, une fois l'organisme sensibilisé.

\corrige{de l'exercice 9}
Face à un microbe, le système immunitaire élimine une réelle menace ; face à un
allergène, il réagit de façon excessive contre une substance qui n'est pourtant pas
dangereuse.

\corrige{de l'exercice 10}
Parce qu'ils bloquent les effets de l'histamine, la substance responsable des
symptômes de l'allergie.

\corrige{de l'exercice 11}
Parce qu'elle agit rapidement pour contrer les effets dangereux (notamment
respiratoires et circulatoires) du choc anaphylactique, qui peut menacer la vie en
quelques minutes.

\corrige{de l'exercice 12}
Lors de la première piqûre, l'organisme se sensibilise sans réaction grave visible ;
lors de la deuxième piqûre, l'organisme, déjà sensibilisé, libère massivement de
l'histamine, ce qui provoque le choc anaphylactique.

\corrige{du sujet 1}
1) Un allergène est une substance inoffensive qui déclenche une réaction excessive du
système immunitaire ; l'allergie est cette réaction excessive et inappropriée.\\
2) Premier contact : sensibilisation, sans symptôme visible. Contact suivant :
libération d'histamine et apparition des symptômes.\\
3) Par exemple : éternuements, démangeaisons, éruption cutanée.\\
4) Par exemple : l'éviction de l'allergène et la prise d'antihistaminiques.

\end{multicols}
\exercicesBanner
